% =========================
% Glosario — Sistemas y Aplicaciones Distribuidas
% =========================
\newglossaryentry{adaptacion}{
	name={Adaptación},
	text={adaptación},
	first={adaptación},
	plural={adaptaciones},
	firstplural={adaptaciones},
	sort={adaptacion},
	description={Capacidad de un sistema o componente para ajustarse a cambios en el entorno, requisitos o condiciones operativas sin perder funcionalidad esencial.},
	symbol={🔄},
	see={reconfiguracion,estabilidad},
	type={general}
}

\newglossaryentry{adaptacioncontrolada}{
	name={Adaptación Controlada},
	text={adaptación controlada},
	first={adaptación controlada},
	plural={adaptaciones controladas},
	firstplural={adaptaciones controladas},
	sort={adaptacioncontrolada},
	description={Proceso de adaptación gobernado por reglas, políticas o mecanismos definidos que limitan y supervisan los cambios realizados en un sistema.},
	symbol={🎛},
	see={adaptacion,reconfiguracion},
	type={general}
}

\newglossaryentry{aplicacion}{
	name={Aplicación},
	text={aplicación},
	first={aplicación},
	plural={aplicaciones},
	firstplural={aplicaciones},
	sort={aplicacion},
	description={Programa informático diseñado para realizar una o varias tareas específicas para el usuario, ejecutándose sobre un sistema computacional determinado.},
	symbol={🧩},
	type=main
}

\newglossaryentry{aplicaciondistribuida}{
	name={Aplicación Distribuida},
	text={aplicación distribuida},
	first={aplicación distribuida},
	plural={aplicaciones distribuidas},
	firstplural={aplicaciones distribuidas},
	sort={aplicaciondistribuida},
	description={Aplicación cuyo procesamiento se reparte entre múltiples nodos interconectados que cooperan mediante mecanismos de comunicación para ofrecer un servicio unificado.},
	symbol={🌐},
	type=main
}

\newglossaryentry{arquitectonico}{
	name={Arquitectónico},
	text={arquitectónico},
	first={arquitectónico},
	plural={arquitectónicos},
	firstplural={arquitectónicos},
	sort={arquitectonico},
	description={Relativo a la arquitectura de un sistema, incluyendo su estructura, componentes, relaciones y principios de diseño.},
	symbol={🏗},
	see={arquitectura,enfoquearquitectonico},
	type={general}
}

\newglossaryentry{arquitectura}{
	name={Arquitectura},
	text={arquitectura},
	first={arquitectura},
	plural={arquitecturas},
	firstplural={arquitecturas},
	sort={arquitectura},
	description={Estructura fundamental de un sistema que define componentes, relaciones, principios de diseño y restricciones técnicas.},
	symbol={🏗},
	type=main
}

\newglossaryentry{arquitecturadistribuida}{
	name={Arquitectura Distribuida},
	text={arquitectura distribuida},
	first={arquitectura distribuida},
	plural={arquitecturas distribuidas},
	firstplural={arquitecturas distribuidas},
	sort={arquitecturadistribuida},
	description={Arquitectura en la que los componentes del sistema se encuentran desplegados en distintos nodos y coordinados mediante redes de comunicación.},
	symbol={🗺},
	type=main
}

\newglossaryentry{asincrona}{
	name={Asíncrona},
	text={asíncrona},
	first={asíncrona},
	plural={asíncronas},
	firstplural={asíncronas},
	sort={asincrona},
	description={Característica de un proceso o mecanismo de comunicación en el que las operaciones no requieren sincronización temporal estricta entre emisor y receptor, permitiendo la ejecución independiente y no bloqueante.},
	symbol={⏳},
	type=main
}

\newglossaryentry{coherencia}{
	name={Coherencia},
	text={coherencia},
	first={coherencia},
	plural={coherencias},
	firstplural={coherencias},
	sort={coherencia},
	description={Grado de consistencia lógica y alineación entre los elementos que conforman un sistema o modelo conceptual.},
	symbol={🧩},
	see={coherenciainterna,arquitectura},
	type={general}
}

\newglossaryentry{coherenciainterna}{
	name={Coherencia Interna},
	text={coherencia interna},
	first={coherencia interna},
	plural={coherencias internas},
	firstplural={coherencias internas},
	sort={coherenciainterna},
	description={Propiedad que asegura que los componentes y decisiones internas de un sistema no presenten contradicciones entre sí.},
	symbol={✔},
	see={coherencia,arquitectura},
	type={general}
}

\newglossaryentry{componente}{
	name={Componente},
	text={componente},
	first={componente},
	plural={componentes},
	firstplural={componentes},
	sort={componente},
	description={Unidad modular de software con responsabilidades bien definidas, capaz de interactuar con otros componentes a través de interfaces.},
	symbol={🔩},
	type=main
}

\newglossaryentry{comunicacion}{
	name={Comunicación},
	text={comunicación},
	first={comunicación},
	plural={comunicaciones},
	firstplural={comunicaciones},
	sort={comunicacion},
	description={Proceso mediante el cual dos o más entidades intercambian información a través de un medio o protocolo definido.},
	symbol={📡},
	type=main
}

\newglossaryentry{comunicacionasincrona}{
	name={Comunicación Asíncrona},
	text={comunicación asíncrona},
	first={comunicación asíncrona},
	plural={comunicaciones asíncronas},
	firstplural={comunicaciones asíncronas},
	sort={comunicacionasincrona},
	description={Modelo de comunicación en el que el emisor no espera una respuesta inmediata del receptor.},
	symbol={⏳},
	type=main
}

\newglossaryentry{comunicacionprocesos}{
	name={Comunicación entre Procesos},
	text={comunicación entre procesos},
	first={comunicación entre procesos},
	plural={comunicaciones entre procesos},
	firstplural={comunicaciones entre procesos},
	sort={comunicacionprocesos},
	description={Mecanismo que permite el intercambio de datos y sincronización entre procesos que se ejecutan de manera concurrente.},
	symbol={🔄},
	type=main
}

\newglossaryentry{comunicacionsincrona}{
	name={Comunicación Síncrona},
	text={comunicación síncrona},
	first={comunicación síncrona},
	plural={comunicaciones síncronas},
	firstplural={comunicaciones síncronas},
	sort={comunicacionsincrona},
	description={Modelo de comunicación en el que el emisor espera la respuesta del receptor antes de continuar su ejecución.},
	symbol={⏱},
	type=main
}

\newglossaryentry{consistenciadatos}{
	name={Consistencia de Datos},
	text={consistencia de datos},
	first={consistencia de datos},
	plural={consistencias de datos},
	firstplural={consistencias de datos},
	sort={consistenciadatos},
	description={Propiedad que garantiza que todos los nodos de un sistema distribuido observan el mismo estado lógico de los datos.},
	symbol={✔},
	type=main
}

\newglossaryentry{disponibilidad}{
	name={Disponibilidad},
	text={disponibilidad},
	first={disponibilidad},
	plural={disponibilidades},
	firstplural={disponibilidades},
	sort={disponibilidad},
	description={Capacidad de un sistema para permanecer operativo y accesible durante un periodo determinado.},
	symbol={🟢},
	type=main
}

\newglossaryentry{enfoque}{
	name={Enfoque},
	text={enfoque},
	first={enfoque},
	plural={enfoques},
	firstplural={enfoques},
	sort={enfoque},
	description={Perspectiva conceptual o metodológica adoptada para analizar, diseñar o resolver un problema determinado.},
	symbol={🔍},
	see={enfoquearquitectonico},
	type={general}
}

\newglossaryentry{enfoquearquitectonico}{
	name={Enfoque Arquitectónico},
	text={enfoque arquitectónico},
	first={enfoque arquitectónico},
	plural={enfoques arquitectónicos},
	firstplural={enfoques arquitectónicos},
	sort={enfoquearquitectonico},
	description={Orientación de diseño que prioriza decisiones estructurales, patrones y estilos arquitectónicos para satisfacer atributos de calidad del sistema.},
	symbol={🗺},
	see={arquitectura,arquitectonico},
	type={general}
}

\newglossaryentry{escalabilidad}{
	name={Escalabilidad},
	text={escalabilidad},
	first={escalabilidad},
	plural={escalabilidades},
	firstplural={escalabilidades},
	sort={escalabilidad},
	description={Capacidad de un sistema para aumentar su rendimiento al incrementar recursos computacionales.},
	symbol={📈},
	type=main
}

\newglossaryentry{escalable}{
	name={Escalable},
	text={escalable},
	first={escalable},
	plural={escalables},
	firstplural={escalables},
	sort={escalable},
	description={Propiedad de un sistema que puede adaptarse eficientemente al crecimiento de la carga de trabajo.},
	symbol={➕},
	type=main
}


\newglossaryentry{estabilidad}{
	name={Estabilidad},
	text={estabilidad},
	first={estabilidad},
	plural={estabilidades},
	firstplural={estabilidades},
	sort={estabilidad},
	description={Capacidad de un sistema para mantener un comportamiento predecible y controlado ante cambios internos o externos.},
	symbol={⚖},
	see={adaptacion,operatividad},
	type={general}
}

\newglossaryentry{interaccion}{
	name={Interacción},
	text={interacción},
	first={interacción},
	plural={interacciones},
	firstplural={interacciones},
	sort={interaccion},
	description={Proceso mediante el cual dos o más entidades intercambian información o coordinan acciones dentro de un sistema.},
	symbol={🤝},
	see={comunicacion,componente},
	type={general}
}

\newglossaryentry{interconectado}{
	name={Interconectado},
	text={interconectado},
	first={interconectado},
	plural={interconectados},
	firstplural={interconectados},
	sort={interconectado},
	description={Estado en el que múltiples nodos o sistemas se encuentran enlazados mediante una red de comunicación.},
	symbol={🔗},
	type=main
}

\newglossaryentry{mecanismo}{
	name={Mecanismo},
	text={mecanismo},
	first={mecanismo},
	plural={mecanismos},
	firstplural={mecanismos},
	sort={mecanismo},
	description={Conjunto de técnicas o procedimientos utilizados para implementar una funcionalidad específica.},
	symbol={⚙},
	type=main
}

\newglossaryentry{mecanismorecuperacion}{
	name={Mecanismo de Recuperación},
	text={mecanismo de recuperación},
	first={mecanismo de recuperación},
	plural={mecanismos de recuperación},
	firstplural={mecanismos de recuperación},
	sort={mecanismorecuperacion},
	description={Estrategia que permite restaurar el funcionamiento normal de un sistema tras una falla.},
	symbol={♻},
	type=main
}

\newglossaryentry{nodo}{
	name={Nodo},
	text={nodo},
	first={nodo},
	plural={nodos},
	firstplural={nodos},
	sort={nodo},
	description={Unidad computacional individual dentro de un sistema distribuido.},
	symbol={🖥},
	type=main
}

\newglossaryentry{operatividad}{
	name={Operatividad},
	text={operatividad},
	first={operatividad},
	plural={operatividades},
	firstplural={operatividades},
	sort={operatividad},
	description={Grado en que un sistema se encuentra en condiciones de funcionamiento efectivo y continuo.},
	symbol={⚙},
	see={disponibilidad,estabilidad},
	type={general}
}

\newglossaryentry{operatividadsistema}{
	name={Operatividad del Sistema},
	text={operatividad del sistema},
	first={operatividad del sistema},
	plural={operatividades del sistema},
	firstplural={operatividades del sistema},
	sort={operatividadsistema},
	description={Estado global que refleja la capacidad del sistema para ejecutar correctamente sus funciones bajo condiciones reales de operación.},
	symbol={🟢},
	see={operatividad,disponibilidad,sistemadistribuido},
	type={general}
}

\newglossaryentry{protocolo}{
	name={Protocolo},
	text={protocolo},
	first={protocolo},
	plural={protocolos},
	firstplural={protocolos},
	sort={protocolo},
	description={Conjunto de reglas que gobiernan la comunicación entre entidades de un sistema.},
	symbol={📜},
	type=main
}

\newglossaryentry{protocoloestandarizado}{
	name={Protocolo Estandarizado},
	text={protocolo estandarizado},
	first={protocolo estandarizado},
	plural={protocolos estandarizados},
	firstplural={protocolos estandarizados},
	sort={protocoloestandarizado},
	description={Protocolo definido y aceptado por organismos de estandarización para asegurar interoperabilidad.},
	symbol={🏛},
	type=main
}

\newglossaryentry{reconfiguracion}{
	name={Reconfiguración},
	text={reconfiguración},
	first={reconfiguración},
	plural={reconfiguraciones},
	firstplural={reconfiguraciones},
	sort={reconfiguracion},
	description={Modificación dinámica de la estructura, configuración o comportamiento de un sistema para responder a cambios operativos.},
	symbol={🔧},
	see={adaptacion,adaptacioncontrolada},
	type={general}
}

\newglossaryentry{recuperacion}{
	name={Recuperación},
	text={recuperación},
	first={recuperación},
	plural={recuperaciones},
	firstplural={recuperaciones},
	sort={recuperacion},
	description={Proceso mediante el cual un sistema vuelve a un estado operativo tras una interrupción.},
	symbol={🔁},
	type=main
}

\newglossaryentry{recursocomputacional}{
	name={Recurso Computacional},
	text={recurso computacional},
	first={recurso computacional},
	plural={recursos computacionales},
	firstplural={recursos computacionales},
	sort={recursocomputacional},
	description={Elemento físico o lógico utilizado para el procesamiento, almacenamiento o comunicación de información.},
	symbol={💾},
	type=main
}

\newglossaryentry{redundancia}{
	name={Redundancia},
	text={redundancia},
	first={redundancia},
	plural={redundancias},
	firstplural={redundancias},
	sort={redundancia},
	description={Duplicación de componentes o datos para aumentar la confiabilidad del sistema.},
	symbol={📦},
	type=main
}

\newglossaryentry{rendimiento}{
	name={Rendimiento},
	text={rendimiento},
	first={rendimiento},
	plural={rendimientos},
	firstplural={rendimientos},
	sort={rendimiento},
	description={Medida de la eficiencia con la que un sistema ejecuta sus operaciones.},
	symbol={🚀},
	type=main
}

\newglossaryentry{seguridad}{
	name={Seguridad},
	text={seguridad},
	first={seguridad},
	plural={seguridades},
	firstplural={seguridades},
	sort={seguridad},
	description={Propiedad de un sistema informático orientada a proteger datos, recursos y servicios frente a accesos no autorizados, alteraciones, interrupciones o divulgación indebida, garantizando confidencialidad, integridad y disponibilidad.},
	symbol={🔐},
	type={general}
}

\newglossaryentry{seguridadistribuida}{
	name={Seguridad Distribuida},
	text={seguridad distribuida},
	first={seguridad distribuida},
	plural={seguridades distribuidas},
	firstplural={seguridades distribuidas},
	sort={seguridaddistribuida},
	description={Enfoque de seguridad aplicado a sistemas distribuidos que protege nodos, comunicaciones y servicios mediante mecanismos coordinados de autenticación, autorización, cifrado, control de acceso y tolerancia a fallos en entornos descentralizados.},
	symbol={🛡},
	type={general}
}

\newglossaryentry{servicio}{
	name={Servicio},
	text={servicio},
	first={servicio},
	plural={servicios},
	firstplural={servicios},
	sort={servicio},
	description={Funcionalidad ofrecida por un componente o sistema a través de una interfaz definida.},
	symbol={🛎},
	type=main
}

\newglossaryentry{sincrona}{
	name={Síncrona},
	text={síncrona},
	first={síncrona},
	plural={síncronas},
	firstplural={síncronas},
	sort={sincrona},
	description={Característica de un proceso que requiere coordinación temporal estricta entre participantes.},
	symbol={⏲},
	type=main
}

\newglossaryentry{sistema}{
	name={Sistema},
	text={sistema},
	first={sistema},
	plural={sistemas},
	firstplural={sistemas},
	sort={sistema},
	description={Conjunto organizado de componentes que interactúan para cumplir un objetivo común.},
	symbol={🧠},
	type=main
}

\newglossaryentry{sistemadistribuido}{
	name={Sistema Distribuido},
	text={sistema distribuido},
	first={sistema distribuido},
	plural={sistemas distribuidos},
	firstplural={sistemas distribuidos},
	sort={sistemadistribuido},
	description={Sistema compuesto por múltiples nodos autónomos que cooperan mediante comunicación para ofrecer servicios.},
	symbol={🌍},
	type=main
}

\newglossaryentry{subyacente}{
	name={Subyacente},
	text={subyacente},
	first={subyacente},
	plural={subyacentes},
	firstplural={subyacentes},
	sort={subyacente},
	description={Característica que describe aquello que se encuentra implícito, oculto o que sirve de base fundamental para el funcionamiento de un sistema, proceso o tecnología.},
	symbol={⬇},
	type={general}
}

\newglossaryentry{tecnologia}{
	name={Tecnología},
	text={tecnología},
	first={tecnología},
	plural={tecnologías},
	firstplural={tecnologías},
	sort={tecnologia},
	description={Conjunto de conocimientos, técnicas, herramientas y métodos aplicados de forma sistemática para el diseño, desarrollo y operación de sistemas informáticos.},
	symbol={🧪},
	type={general}
}

\newglossaryentry{tecnologiasubyacente}{
	name={Tecnología Subyacente},
	text={tecnología subyacente},
	first={tecnología subyacente},
	plural={tecnologías subyacentes},
	firstplural={tecnologías subyacentes},
	sort={tecnologiasubyacente},
	description={Tecnología fundamental que soporta y habilita el funcionamiento de un sistema o aplicación, incluyendo plataformas, protocolos, infraestructuras y mecanismos sobre los cuales se construyen servicios de mayor nivel.},
	symbol={🧱},
	type={general}
}

\newglossaryentry{toleranciafallos}{
	name={Tolerancia a Fallos},
	text={tolerancia a fallos},
	first={tolerancia a fallos},
	plural={tolerancias a fallos},
	firstplural={tolerancias a fallos},
	sort={toleranciafallos},
	description={Capacidad de un sistema para continuar operando correctamente ante fallas parciales.},
	symbol={🛡},
	type=main
}

\newglossaryentry{tolerantefallos}{
	name={Tolerante a Fallos},
	text={tolerante a fallos},
	first={tolerante a fallos},
	plural={tolerantes a fallos},
	firstplural={tolerantes a fallos},
	sort={tolerantefallos},
	description={Propiedad de un sistema diseñado para resistir fallas sin interrumpir su funcionamiento.},
	symbol={🧱},
	type=main
}
